
\documentclass{article}
\usepackage{colortbl}
\usepackage{makecell}
\usepackage{multirow}
\usepackage{supertabular}

\begin{document}

\newcounter{utterance}

\centering \large Interaction Transcript for game `cladder', experiment `full\_v1.5\_default', episode 6126 with Qwen3.5{-}4B.
\vspace{24pt}

{ \footnotesize  \setcounter{utterance}{1}
\setlength{\tabcolsep}{0pt}
\begin{supertabular}{c@{$\;$}|p{.15\linewidth}@{}p{.15\linewidth}p{.15\linewidth}p{.15\linewidth}p{.15\linewidth}p{.15\linewidth}}
    \# & \multicolumn{2}{c}{Player} && \multicolumn{2}{c}{Game Master} \\
    \hline

    \theutterance \stepcounter{utterance}  
    & & & \multicolumn{4}{p{0.6\linewidth}}{
        \cellcolor[rgb]{0.9,0.9,0.9}{
            \makecell[{{p{\linewidth}}}]{
                \texttt{\tiny{[P1$\langle$GM]}}
                \texttt{You are an expert in causal inference. The following question is not a typical commonsense query, but rather a meticulously designed question created by a professor specializing in causal inference, intended to assess the students' mastery of the course content.} \\
\texttt{Imagine a self{-}contained, hypothetical world with only the following conditions, and without any unmentioned factors or causal relationships: Smoking has a direct effect on effort and lactose intolerance. Effort has a direct effect on lactose intolerance. We know that smoking causes being lazy. smoking and being hard{-}working causes being lactose intolerant. Would the student is lactose intolerant if nonsmoking instead of smoking?} \\
\texttt{Start your answer with 'Yes' or 'No', followed by additional reasoning or evidence to support your explanation.} \\
            }
        }
    }
     \\ \\

    \theutterance \stepcounter{utterance}  
    & \multicolumn{4}{p{0.6\linewidth}}{
        \cellcolor[rgb]{0.561,0.737,0.561}{
            \makecell[{{p{\linewidth}}}]{
                \texttt{\tiny{[P1$\rangle$GM]}}
                \texttt{No} \\
\\ 
\texttt{**Reasoning:**} \\
\\ 
\texttt{To determine the outcome, we must trace the causal pathways defined in this hypothetical world, strictly adhering to the provided conditions and ignoring external real{-}world knowledge (such as the actual biological link between smoking and lactose intolerance).} \\
\\ 
\texttt{1.  **Analyze the Causal Chain:**} \\
\texttt{    *   **Condition A:** Smoking \$\textbackslash{}rightarrow\$ Being lazy.} \\
\texttt{    *   **Condition B:** Being lazy \$\textbackslash{}rightarrow\$ Effort (implied: being lazy reduces effort).} \\
\texttt{    *   **Condition C:** Effort \$\textbackslash{}rightarrow\$ Lactose intolerance.} \\
\texttt{    *   **Condition D:** Smoking \$\textbackslash{}rightarrow\$ Lactose intolerance (direct effect).} \\
\texttt{    *   **Condition E:** Being hard{-}working \$\textbackslash{}rightarrow\$ Lactose intolerance (direct effect).} \\
\\ 
\texttt{2.  **Evaluate the Scenario:**} \\
\texttt{    *   The subject is **nonsmoking**.} \\
\texttt{    *   Since smoking causes being lazy, **nonsmoking** implies the absence of the cause for laziness. In this closed{-}world assumption, if the cause (smoking) is removed, the effect (laziness) does not occur. Therefore, the subject is **not lazy**.} \\
\texttt{    *   If the subject is not lazy, they are **hard{-}working** (assuming a binary state of effort where not lazy equals hard{-}working, or simply that the specific causal chain for low effort is broken).} \\
\texttt{    *   According to **Condition E**, being hard{-}working has a direct effect on causing lactose intolerance.} \\
\\ 
\texttt{3.  **Conclusion:**} \\
\texttt{    Even though the subject is not smoking, the removal of smoking triggers the reverse causal chain: Nonsmoking \$\textbackslash{}rightarrow\$ Not Lazy \$\textbackslash{}rightarrow\$ Hard{-}working \$\textbackslash{}rightarrow\$ **Lactose Intolerance**.} \\
\\ 
\texttt{Therefore, the student would still be lactose intolerant.} \\
            }
        }
    }
    & & \\ \\

    \theutterance \stepcounter{utterance}  
    & & & \multicolumn{2}{p{0.3\linewidth}}{
        \cellcolor[rgb]{0.9,0.9,0.9}{
            \makecell[{{p{\linewidth}}}]{
                \texttt{\tiny{[GM$|$GM]}}
                \texttt{No} \\
            }
        }
    }
    & & \\ \\

    \theutterance \stepcounter{utterance}  
    & & & \multicolumn{2}{p{0.3\linewidth}}{
        \cellcolor[rgb]{0.9,0.9,0.9}{
            \makecell[{{p{\linewidth}}}]{
                \texttt{\tiny{[GM$|$GM]}}
                \texttt{no} \\
            }
        }
    }
    & & \\ \\

    \theutterance \stepcounter{utterance}  
    & & & \multicolumn{2}{p{0.3\linewidth}}{
        \cellcolor[rgb]{0.9,0.9,0.9}{
            \makecell[{{p{\linewidth}}}]{
                \texttt{\tiny{[GM$|$GM]}}
                \texttt{game\_result = WIN} \\
            }
        }
    }
    & & \\ \\

\end{supertabular}
}

\end{document}
